% GENERATED FILE. Edit canonical lessons, then run npm run generate:books.
% canonical-source-hash: fnv1a64:664c5fbe18c4e46e
% canonical-lessons: GE-C262-zunaechst, GE-C262-anschliessend, GE-C262-daraufhin, GE-C262-zuletzt, GE-C262-reihenfolge

\chapter{Zunächst und anschließend}
\label{ch:ge-zunachst-und-anschlie-end}
\hlchaptermodality{262}

\begin{chapteropening}
\textbf{By the end of this chapter:} \emph{I can put the steps of a story in order: at first, afterwards, as a result, last of all and the order of things.}

Complete the last lesson of chapter 262: I can put the steps of a story in order: at first, afterwards, as a result, last of all and the order of things.
\end{chapteropening}

\section[zunächst]{zunächst — at first, initially; for now}

\label{lesson:GE-C262-zunaechst}



\begin{quote}
Before the new one: say the German for at least, then the German for during.
\end{quote}

\subsection*{You'll want to know: zunächst}
\textbf{zunächst} — \textquotedblleft{}at first, initially; for now\textquotedblright{}.

\begin{quote}
Zunächst war alles ruhig.
\end{quote}

\begin{quote}
\emph{At first everything was quiet.}
\end{quote}

A little more formal than \emph{zuerst}. It can also mean 'for now': \emph{zunächst bleibe ich zu Hause}.

\begin{cousinweb}
\emph{zunächst} is built on \emph{zu} \textquotedblleft{}to\textquotedblright{} and \emph{nächst} \textquotedblleft{}nearest\textquotedblright{}, the superlative of \emph{nah} \textquotedblleft{}near\textquotedblright{}, which is inherited from Old High German \emph{nāh} \textquotedblleft{}close, near\textquotedblright{} (EtymDB): the nearest step comes first.
\end{cousinweb}

\begin{grammarlens}[title={What you've built}]
One more word for telling, explaining and describing at B1.
\end{grammarlens}

\subsection*{Your turn}
\begin{itemize}
\raggedright
  \item \emph{Say it:} \emph{zunächst}
  \item \emph{Say it:} \emph{zunächst}, once more
  \item \emph{Say it:} a sentence of your own with \emph{zunächst}
  \item \emph{Recall:} say the German for at least, then the German for during, then say \emph{zunächst} again
\end{itemize}

\begin{tcolorbox}[breakable,colback=teal!4,colframe=teal!35!black,title={Before you move on}]
What does zunächst mean? (\textquotedblleft{}At first, initially; for now\textquotedblright{}.) Say it once more.
\end{tcolorbox}
\section[anschließend]{anschließend — afterwards, subsequently}

\label{lesson:GE-C262-anschliessend}



\begin{quote}
Before the new one: say the German for during, then the German for at first.
\end{quote}

\subsection*{You'll want to know: anschließend}
\textbf{anschließend} — \textquotedblleft{}afterwards, subsequently\textquotedblright{}.

\begin{quote}
Wir waren im Kino, und anschließend haben wir gegessen.
\end{quote}

\begin{quote}
\emph{We were at the cinema, and afterwards we ate.}
\end{quote}

At the start of a clause the verb comes second: \emph{Anschließend gingen wir nach Hause}.

\begin{cousinweb}
\emph{anschließend} is the present participle of \emph{anschließen} \textquotedblleft{}to join on\textquotedblright{}, built on \emph{schließen} \textquotedblleft{}to close\textquotedblright{}, which is inherited from Old High German \emph{sliozan} \textquotedblleft{}to close, to lock\textquotedblright{} (EtymDB): what is joined on next.
\end{cousinweb}

\begin{grammarlens}[title={What you've built}]
One more word for telling, explaining and describing at B1.
\end{grammarlens}

\subsection*{Your turn}
\begin{itemize}
\raggedright
  \item \emph{Say it:} \emph{anschließend}
  \item \emph{Say it:} \emph{anschließend}, once more
  \item \emph{Say it:} a sentence of your own with \emph{anschließend}
  \item \emph{Recall:} say the German for during, then the German for at first, then say \emph{anschließend} again
\end{itemize}

\begin{tcolorbox}[breakable,colback=teal!4,colframe=teal!35!black,title={Before you move on}]
What does anschließend mean? (\textquotedblleft{}Afterwards, subsequently\textquotedblright{}.) Say it once more.
\end{tcolorbox}
\section[daraufhin]{daraufhin — as a result, thereupon}

\label{lesson:GE-C262-daraufhin}



\begin{quote}
Before the new one: say the German for at first, then the German for afterwards.
\end{quote}

\subsection*{You'll want to know: daraufhin}
\textbf{daraufhin} — \textquotedblleft{}as a result, thereupon\textquotedblright{}.

\begin{quote}
Der Bus kam nicht. Daraufhin nahmen wir ein Taxi.
\end{quote}

\begin{quote}
\emph{The bus did not come. As a result, we took a taxi.}
\end{quote}

It says that one thing led to the next. The verb comes second: \emph{Daraufhin rief sie an}.

\begin{cousinweb}
\emph{daraufhin} joins \emph{darauf} \textquotedblleft{}on that, after that\textquotedblright{} and \emph{hin} \textquotedblleft{}towards, away\textquotedblright{}, all inherited Germanic pieces; EtymDB lists no parent for the whole, which is a modern compound.
\end{cousinweb}

\begin{grammarlens}[title={What you've built}]
One more word for telling, explaining and describing at B1.
\end{grammarlens}

\subsection*{Your turn}
\begin{itemize}
\raggedright
  \item \emph{Say it:} \emph{daraufhin}
  \item \emph{Say it:} \emph{daraufhin}, once more
  \item \emph{Say it:} a sentence of your own with \emph{daraufhin}
  \item \emph{Recall:} say the German for at first, then the German for afterwards, then say \emph{daraufhin} again
\end{itemize}

\begin{tcolorbox}[breakable,colback=teal!4,colframe=teal!35!black,title={Before you move on}]
What does daraufhin mean? (\textquotedblleft{}As a result, thereupon\textquotedblright{}.) Say it once more.
\end{tcolorbox}
\section[zuletzt]{zuletzt — last of all, finally; the last time}

\label{lesson:GE-C262-zuletzt}



\begin{quote}
Before the new one: say the German for afterwards, then the German for as a result.
\end{quote}

\subsection*{You'll want to know: zuletzt}
\textbf{zuletzt} — \textquotedblleft{}last of all, finally; the last time\textquotedblright{}.

\begin{quote}
Zuletzt hat er die Tür zugemacht.
\end{quote}

\begin{quote}
\emph{Last of all, he shut the door.}
\end{quote}

Also 'the last time': \emph{Wann hast du sie zuletzt gesehen?}

\begin{cousinweb}
\emph{zuletzt} joins \emph{zu} and \emph{letzt} \textquotedblleft{}last\textquotedblright{}. \emph{Letzt} is usually traced to Old High German \emph{lezzisto}, the superlative of \emph{laz} \textquotedblleft{}slow, weary\textquotedblright{}, a cousin of English \emph{late}; EtymDB lists no parent.
\end{cousinweb}

\begin{grammarlens}[title={What you've built}]
One more word for telling, explaining and describing at B1.
\end{grammarlens}

\subsection*{Your turn}
\begin{itemize}
\raggedright
  \item \emph{Say it:} \emph{zuletzt}
  \item \emph{Say it:} \emph{zuletzt}, once more
  \item \emph{Say it:} a sentence of your own with \emph{zuletzt}
  \item \emph{Recall:} say the German for afterwards, then the German for as a result, then say \emph{zuletzt} again
\end{itemize}

\begin{tcolorbox}[breakable,colback=teal!4,colframe=teal!35!black,title={Before you move on}]
What does zuletzt mean? (\textquotedblleft{}Last of all, finally; the last time\textquotedblright{}.) Say it once more.
\end{tcolorbox}
\section[die Reihenfolge]{die Reihenfolge — the order, the sequence}

\label{lesson:GE-C262-reihenfolge}



\begin{quote}
Before the new one: say the German for as a result, then the German for last of all.
\end{quote}

\subsection*{You'll want to know: die Reihenfolge}
\textbf{die Reihenfolge} — \textquotedblleft{}the order, the sequence\textquotedblright{}.

\begin{quote}
Die Reihenfolge der Bilder war falsch.
\end{quote}

\begin{quote}
\emph{The order of the pictures was wrong.}
\end{quote}

Plural: \emph{die Reihenfolgen}. \emph{In der richtigen Reihenfolge} — in the right order.

\begin{cousinweb}
\emph{Reihenfolge} joins \emph{Reihe} \textquotedblleft{}row\textquotedblright{}, from Middle High German \emph{rīhe} \textquotedblleft{}row, line\textquotedblright{}, and \emph{Folge} \textquotedblleft{}what follows\textquotedblright{}, from Middle High German \emph{volge} (EtymDB). Both are inherited Germanic words.
\end{cousinweb}

\begin{grammarlens}[title={What you've built}]
One more word for telling, explaining and describing at B1.
\end{grammarlens}

\subsection*{Your turn}
\begin{itemize}
\raggedright
  \item \emph{Say it:} \emph{die Reihenfolge}
  \item \emph{Say it:} \emph{die Reihenfolge}, once more
  \item \emph{Say it:} a sentence of your own with \emph{die Reihenfolge}
  \item \emph{Recall:} say the German for as a result, then the German for last of all, then say \emph{die Reihenfolge} again
\end{itemize}

\begin{tcolorbox}[breakable,colback=teal!4,colframe=teal!35!black,title={Before you move on}]
What does die Reihenfolge mean? (\textquotedblleft{}The order, the sequence\textquotedblright{}.) Say it once more.
\end{tcolorbox}
